\documentclass[main.tex]{subfiles}

\begin{document}

\section*{Dicussion}

Our simulation study demonstrates that introducing DRT in Cape Town can effectively reduce traffic congestion, with the extent of relief depending on fleet size and pricing strategies. Across 60 scenarios, DRT consistently decreased congestion, with the most effective configurations achieving up to 28.2\,\% citywide reductions. Notably, DRT not only reduced the number of private cars on the road but also improved traffic flow efficiency, as evidenced by congestion reductions often exceeding the decrease in vehicle numbers. This suggests that the benefits of DRT extend beyond mere vehicle substitution, likely through optimized routing and increased vehicle occupancy via shared rides. 

The analysis of mode choice shifts reveals that DRT does not significantly cannibalize public transport. Instead, it often leads to a net increase in public transport usage by facilitating multi-modal trips and attracting former car users. This finding is crucial for Cape Town, where public transport is overburdened, and reducing car dependency is a priority. DRT could add capacity that public transport does not provide, and at the same time substitute private cars. However, the introduction of DRT also adds to private motorized mileage, which could have environmental implications, though these are mitigated by the efficiency of shared rides and the potential for electric DRT fleets. 

Ride-hailing services such as Uber and Bolt, which hold licenses to operate in Cape Town, were not explicitly modeled in this study to focus on DRT as a shared, efficiency-oriented alternative. These services typically feature pricing skewed toward car costs (often 100\,\% or more of variable car fees), unshared rides, and significant dead mileage, which can worsen congestion rather than alleviate it \citep{Erhardt2019-rs} \citep{Tirachini2020-sr}. By positioning DRT at lower price points (e.g., 45.45–75\,\% of car km fees), our scenarios demonstrate how DRT could compete with ride-hailing by attracting car users to a more sustainable option, potentially reducing overall vehicle kilometers and congestion. Future iterations of the model could incorporate ride-hailing to simulate competitive dynamics, as their inclusion might influence optimal DRT fleet sizes (e.g., requiring larger fleets to offset ride-hailing's market share) and pricing (e.g., to undercut unshared services). 

Geographically, the buffer areas around DRT service zones experienced the largest congestion reductions, indicating that DRT can alleviate traffic bottlenecks at the edges of high-car-demand areas. Within service areas, congestion occasionally increased in expensive DRT scenarios, underscoring the importance of competitive pricing to ensure adoption and effectiveness. The scenario analysis highlights DRT’s potential to significantly reduce congestion in critical areas like the CBD, where traffic issues are most acute. This targeted impact suggests that strategically deploying DRT in high-traffic zones can yield substantial benefits for urban mobility.

While our study provides valuable insights into DRT's potential in Cape Town, several limitations should be acknowledged. First, while pricing assumptions draw on empirical references like MOIA's fare structure, they may overlook dynamic costing factors such as fluctuating subsidies, maintenance expenses, or inflation-adjusted long-term viability in a South African context. Second, our approach assumes higher adoption in predefined upper-middle-income areas, but it may not fully capture real-world behavioural barriers to DRT uptake. Residents in Cape Town often exhibit a strong preference for private cars even in congested conditions, driven by distrust in shared systems \citep{Konig2018-pe}. These psychological factors, well-documented in the mobility literature and widely recognised by local experts, could result in smaller mode shifts than simulated. The presented congestion reductions are therefore conditional on successful trust-building measures (e.g., pilot programmes, safety certifications, or communication campaigns) that go beyond pure price and fleet-size incentives. Third, we did not analyze the effect of DRT as an efficient "first-mile/last-mile" connector. In poorly designed public transportation systems and rapid speed of urbanization, DRT might be enhancing the overall reach and utility of existing high-capacity mass transit systems and thus drawing more commuters away from private vehicles, ultimately moving more people with less road space.

These limitations highlight opportunities for future research, including real-world trials to validate costing and availability, and expanded models incorporating socioeconomic incentives to ensure inclusive DRT implementation. Furthermore, there is a clear need for choice modelling studies to rigorously test users' willingness to share and quantify the trade-offs between cost, time, and privacy. Such research, alongside effective strategies to reduce sharing anxiety—for example through trust-building pilot programmes, safety certifications, transparent vehicle-tracking features, or community engagement campaigns—could significantly lower psychological barriers to DRT adoption and thereby realise even greater mode shifts and congestion relief than those simulated. Future research should also explore the integration of DRT with existing public transport systems, the optimization of fleet operations to minimize dead mileage, and the potential of electric DRT vehicles to enhance environmental outcomes. Additionally, policy incentives may be necessary to encourage DRT adoption, particularly in scenarios where DRT travel times exceed those of private cars, such as DRT-designated driving lanes or discounted or free parking for DRT vehicles at transport hubs.


\end{document}